\SourceLeaf{442}%
war, so ist
\[
\begin{aligned}
\chi_1&=\chi+\operatorname{Area}L'L'_1P'_1P'\\
\psi'-\varphi'+A&=LL'_1L'L+\chi_1-\chi;
\end{aligned}
\]
allein
\[
\begin{aligned}
\varphi'&=\chi+B\\
\psi'&=\chi_1+B_1,
\end{aligned}
\]
also
\[
B_1-B+A=LL'_1L'L.
\]
Die Winkel des Dreiecks $LL'L'_1$ sind offenbar
\[
A,\qquad 180^{0}-B,\qquad B_1,
\]
also ihre Summe
\[
180^{0}+LL'_1L'L.
\]

Der Beweis wird in der Form einiger Modification und Erläuterung bedürfen, wenn der Punkt $(3)$ innerhalb des Dreiecks fällt. Allgemein aber schliessen wir:
\begin{quote}
»Die Summe der drei Winkel eines Dreiecks, welches auf einer beliebigen krummen Fläche durch kürzeste Linien gebildet wird, ist gleich der Summe von $180^{0}$ und dem Inhalt des Dreiecks auf der Hülfskugel, dessen Begrenzung durch die Punkte $L$ gebildet wird, welche den Punkten in der Begrenzung jenes Dreiecks entsprechen, und zwar so, dass der gedachte Inhalt des Dreiecks als positiv oder negativ anzusehen ist, je nachdem es von seiner Begrenzung in demselben Sinn umgeben wird wie die Figur oder im entgegengesetzten«.
\end{quote}
\noindent\hspace*{1.2em}Man schliesst hieraus ferner leicht, dass die Summe aller Winkel eines Polygons von $n$ Seiten, die auf der krummen Fläche kürzeste Linien sind, [gleich] der Summe von $(n-2)180^{0}+$ Inhalt des Polygons auf der Kugelfläche etc.

\begin{center}16.\end{center}
\noindent\hspace*{1.2em}Wenn eine krumme Fläche vollkommen auf die andere abgewickelt werden kann, so behalten dabei offenbar alle Linien auf der ersten Fläche bei der Übertragung auf einander ihre Grösse, eben so wie die Winkel, die durch das
